% !TeX program = xelatex
\documentclass[journal,onecolumn]{IEEEtran}
\usepackage{iftex}
\ifPDFTeX
  \errmessage{This document requires XeLaTeX. In Overleaf: Menu -> Compiler -> XeLaTeX}
\fi
\usepackage{geometry}
\geometry{margin=2.3cm}
\usepackage{amsmath,amssymb,amsthm,mathtools}
\usepackage{booktabs}
\usepackage{hyperref}
\hypersetup{hidelinks}
\usepackage{fontspec}
\usepackage{xeCJK}
\setmainfont{Times New Roman}
\setCJKmainfont{FandolSong-Regular}
\linespread{1.2}

\title{论文提纲v2}
\author{匿名作者}
\date{}

\begin{document}
\maketitle

\begin{abstract}
本文为论文提纲排版稿，面向可审计半自动全景布局标注流程，覆盖研究问题、方法、实验设置、报告口径与预期贡献。
\end{abstract}

\begin{IEEEkeywords}
Semi-automatic annotation, Reliability, Routing, OOD, Auditable protocol
\end{IEEEkeywords}
\section{研究问题}

本研究关注全景室内场景中半自动角点标注的质量保障问题。在无法对所有样本进行全量精标的现实约束下，如何系统性地评估半自动标注的效率、质量与可靠性，并利用场景特异可靠度支持任务路由。

技术主线：采用可审计三段式规程 Pilot→PreScreen→Main；主实验前锁定阈值与规则（MDE、$w_u$ 映射、关键超参数）；在 non-IID/OOD 条件下以 meta-label 与 OOD 风险分数驱动预筛选与路由；输出画像、矩阵与反例库等过程证据。
\subsection{研究问题 RQ}

本文围绕三个核心研究问题展开，分别对应效率、质量与路由三个维度。

RQ1 效率评估: 半自动初始化是否降低真实活跃标注时间？下降的比例是多少？
主终点：Manual vs Semi-Auto 的 $active_time$ 中位数差异（秒）及 95\% CI。$active_time$ 的操作定义、采集实现与聚合口径见 §2.X。
风险控制：分析 time 与 edit 幅度的关系，排除时间缩短仅源于随意修改
RQ2 质量与可靠性: 在部分复核约束下，半自动流程是否改变标注一致性与可靠性分布，能否系统性暴露反例类型？
主终点：$Calibration_manual$ vs $Calibration_semi$ 的 $\mathrm{IAA}_t$ 分布差异 K-S 检验
反例类型：异常编辑、结构性错误如跨门扩张、配对失败、角点漂移、门控失败
可复现筛选规则：低 $\mathrm{IAA}_t<0.7$ + 改动模式异常 + 门控失败；最终反例需专家复核

RQ2b: 加权共识作为干预手段，是否能在不引入循环论证的前提下提升共识稳定性？
定位：加权共识是核心算法输出与管理手段，但不作为唯一评价目标；评价目标优先使用预筛选测试样本的专家参考标准，或未加权一致性指标。
评估：在预注册的参考子集上，对比加权与未加权的标签一致性、反例检出敏感性与稳定性。
RQ3 路由与优化: 在固定标注预算下，能否通过预筛选、工人画像与场景或分布感知路由，在 non-IID 与 OOD 场景中维持一致性并降低失败率？
场景定义：以 difficulty 与 $model_issue$ 的共识标签作为场景代理，优先采用少数高频粗粒度场景 s，避免标签组合导致场景爆炸。
场景特异可靠度：
\[r_u^{\left(s\right)}=\mathrm{median}_{t\in T_s}\,\mathrm{IoU}\left(A_{t,u},C_t^{\left(-u\right)}\right)\]
路由策略：Random，Global 按全局 $r_u$，Stratified 按 $r_u^{\left(s\right)}$
评估协议：离线模拟，在每任务已有的 k 份真实标注中，模拟策略分配任务并取实际提交标注作为输出，以避免因使用真实标签导致的信息泄露与乐观估计。
指标：一致性代理提升、active time、门控失败率

\section{方法}

本章建立改动幅度、质量、可靠性三类指标体系，并明确各指标的适用范围与边界。核心原则如下：改动幅度反映工作量而非正确性；质量与可靠性分开统计；所有剔除与失败均记录原因以保证可追溯性。

\subsection{总体框架：预筛选、画像、审计、路由}

本研究采用可复现的预筛选与工人画像框架。考虑到小样本条件下边做边收敛的强化学习闭环缺乏可审计性，本文转而强调过程性证据与规程透明性。具体采用 Pilot→PreScreen→Main Experiment 的三段式流程，先锁定 MDE、$w_u$ 映射、超参数等阈值与规则，再进入主实验，从而保证可复现并降低事后调参空间。

整体流程分为四个阶段：
(1) 预筛选（Pre-screen）：用小规模盲测任务集建立工人基线，剔除明显不合格或不稳定的工人，形成可用工人池。
(2) 校准（Calibration）：在 $Calibration_manual$ 上估计全局可靠度 $r_u$ 与场景特异可靠度 $r_u^{\left(s\right)}$，并输出置信区间。
(3) 在线路由（Routing）：结合场景信号与 OOD 风险触发器，在预算约束下将高风险任务优先分配给更稳定的工人，并采用序贯冗余与停止准则。
(4) 审计输出（Audit）：输出工人画像散点图（可靠度 vs 异常度）与类型分布、Worker×Scene 矩阵、门控失败原因分布、反例案例库，形成可追溯证据链。

\subsection{标注对象与元信息（Meta-Label as Decision Byproduct）}

标注以角点为主，不强制要求画墙面。每次提交时需同步记录三类信息：Scope、Difficulty 与 Model Issue，这三类元标签作为标注决策的副产品，同时服务于后续的场景分层、工人画像与路由策略。
\subsubsection{Scope 单选，必填，决定是否进入主指标}
In-scope: 仅标注相机所在房间
规则：门框与墙垛清晰时，边界止于门框与墙垛处，不跨门洞
包含：天花内凹与格栅不影响外边界
Out-of-scope（共 4 类）:
几何假设不成立：非 Manhattan 或弧形墙
边界不可判定：开放式或无可靠停止点
错层或多平面：楼梯井或分层天花
证据不足：遮挡或画质差导致无法判断
统计处理：标记为 In-scope 的样本进入主指标口径 I；标记为 OOS 的样本单独报告其比例、子类分布与判定一致性；当同一图像存在标注者分别选择 In-scope 与 OOS 的混合判定时，进行人工复核。
\subsubsection{Difficulty 多选}
用户界面要求：在每次提交时，针对每个 difficulty tag 必须显式选择是/否，避免隐含缺失。系统错误或字段漏传时记录为 NA，NA 必须单独统计和报告以供审计。
数学建模：每个 difficulty tag k 对应二元变量 $z_{t,u,k}\in\left\{0,1\right\}$，其中 1 表示该困难存在，0 表示该困难不存在。NA 仅用于字段缺失或系统错误，单独记录与报告。
Tag 集合（共 6个）:
非常简单 trivial：无明显困难，快速完成
遮挡明显 occlusion：家具或人遮挡墙角
纹理弱 low_texture：纯色墙或光照均匀
拼接缝 seam：全景拼接缝导致错位
反光与玻璃 reflection：镜面伪边界
画质差或遮罩 low_quality：模糊或黑边影响对齐

作用：上述 difficulty 标签用于场景分层、路由信号与分布监控。
\subsubsection{Model Issue 多选，仅在半自动条件下收集}
用户界面要求：在半自动修改模型输出时，针对每个 $model_issue$的 tag 必须显式选择是/否，以避免隐含缺失。系统错误或字段漏传时记录为 NA，NA 仅用于系统/提交缺失并需单独统计与报告以供审计。
数学建模：每个 $model_issue$ tag k 对应二元变量$z_{t,u,k}\in\left\{0,1\right\}$，其中 1 表示该模型失效模式存在，0 表示不存在。NA 仅用于字段缺失或系统错误，并且应单独记录与报告。
Tag 集合（共 8个）:
模型标注质量好 acceptable：无需显著修改
跨门扩张 overextend_adjacent：包含相邻空间
漏标 underextend：只标了局部
过度解析 over_parsing：误判柱子 / 家具为拐角
角点错位 corner_drift：整体偏移
角点重复 corner_duplicate：同一拐角多个点
拓扑崩溃 topology_failure：配对 / 闭合失败
预测失效 fail：需从零重画
两者区别：
Difficulty：场景的客观特征，人工标注同样存在困难
$Model_issue$：模型特有的失效模式，可通过人工修正
\subsection{修改幅度与工作量（半自动标注组）}

以下指标用于量化标注者对模型初始化结果的修改程度。设任务 t 的模型初始化结果为$P_t$，标注者 u 的提交结果为 $A_{t,u}$。
\subsubsection{图像空间 Layout Polygon IoU 改动幅度}
以图像空间 polygon 的交并比衡量从初始化到最终提交的形状改动幅度
\[\mathrm{IoU}\left(X,Y\right)=\frac{\left|X\bigcap Y\right|}{\left|X\cup Y\right|}\]
其中 X,Y 为在全景图像空间上的 layout 投影多边形。当前代码将角点集合先构造成多边形后再计算 IoU，其中角点多边形采用凸包（convex hull）以确保多边形有效。
凸包会将凹形结构外包为凸形，可能降低对跨门扩张或凹结构错误的敏感度，因此该 IoU 在论文中作为改动幅度或一致性的辅助证据，主结论优先依赖更几何一致的指标如 BoundaryRMSE 与标准 layout 2D/3D IoU。
在半自动初始化条件下，本指标默认比较模型初始化与标注者修订后的差异：
\[\mathrm{IoU}_{\mathrm{edit}}\left(t,u\right)=\mathrm{IoU}\left(P_t,A_{t,u}\right)\]
$\mathrm{IoU}_{\mathrm{edit}}$ 越低通常表示改动越大，但不代表改动必然正确，仅用于衡量工作量或改动幅度。
\subsubsection{边界 RMSE}
当角点数量不一致或者配对不稳定时，点对点 RMSE 易失真，所以采用角点生成的上 / 下边界曲线 $y_{ceil}\left(x\right)$ , $y_{floor}\left(x\right)$在离散网格上的误差作为几何改动幅度指标
\[\mathrm{BoundaryRMSE}\left(X,Y\right)=\sqrt{\frac{1}{N}\sum_{i=1}^{N}\left[\left(\Delta y_{ceil}\left(x_i\right)\right)^2+\left(\Delta y_{floor}\left(x_i\right)\right)^2\right]}\]

\subsubsection{角点匹配 RMSE 辅助指标}
该指标仅在角点配对成功且配对覆盖率足够高的情况下启用
预测角点的集合为{$p_1,p_2,\ldots.,p_n},$标注角点集合为{$g_1,g_2,\ldots.,g_m},$两点间欧式距离为
\[d\left(p,g\right)=\sqrt{\left[\left(x_p-x_g\right)^2+\left(y_p-y_g\right)^2\right]}\]
用匈牙利算法在两集合间求最小总代价的匹配π,匹配对上计算:
\[\mathrm{RMSE}_{\mathrm{match}}\left(t,u\right)=\sqrt{\frac{1}{N}\sum_{i=1}^{N}{d\left(p_i,g_{\pi\left(i\right)}\right)^2}}\]
其中 $N=\min{\left(n,m\right)}$，取预测角点与标注角点中数量较少的那个。
2.X 活跃标注时间定义（Active Time）

为明确 RQ1 的主终点度量，特界定活跃标注时间之操作定义与聚合方式。
定义与采集条件：
$active_time$ 指标注者在标注任务页面上的交互活跃时间，而非页面打开挂钟时间。计时器每秒递增 1，当且仅当（1）页面处于可见状态；（2）距最近用户交互（鼠标、键盘、滚轮）不超过 $IDLE_THRESHOLD$ 秒。当页面不可见超过 $PAGE_HIDDEN_THRESHOLD$ 秒应视为中断，需新的交互事件方可恢复计时。阈值在预注册中锁定为 $IDLE_THRESHOLD$ = 15 s、$PAGE_HIDDEN_THRESHOLD$ = 6 s，实验期间不做事后调整。
聚合协议：
按 session 层级分段记录，其中 session 定义为同一浏览器标签页内的连续交互窗口，以 sessionStorage 的唯一标识追踪。
Per-(task, annotator, session)：取该 session 内的最大累积值，用于消除同一 session 内的重复上报。
Per-(task, annotator)：对跨 session 的多个分段累加，覆盖标注者在不同时段或标签页对同一任务的重复访问。
Per-task：取参与标注的各标注者中的最大值，用于任务级度量。
\subsection{多选标签的共识计算（Per-Tag Binary Consensus）}

由于 Difficulty 与 Model Issue 均为多选标签，不存在单一多数类，因此需要专门的共识计算机制。

\subsubsection{基本方案}
采用 Per-tag 二元共识作为主口径；加权共识作为可选干预，见2.3.2节，通过预筛选真值锚点与预注册协议规避循环依赖。
共识计算 对每个 tag k 独立：
记号约定：为简洁起见，下述 $c_{t,k}$、$\mathrm{conf}_{t,k}$、{$\mathrm{disagree}}_{t,k}$ 中任务下标 t 有时省略；实际实现均以任务 t 与 tag k 为索引。
输入：任务 t 的 m个worker 标注$z_{t,1,k},z_{t,2,k},\ldots,z_{t,m,k}$
输出：共识 $c_k$，置信度 $\mathrm{conf}_k$，分歧度 {$\mathrm{disagree}}_k$
统计勾选人数：$n_{selected}=\sum_{u=1}^{m}\mathbb{I}\left[z_{t,u,k}=1\right]$
统计有效人数：$n_{total}=\sum_{u=1}^{m}\mathbb{I}\left[z_{t,u,k}\in\{0,1\}\right]$（理想情况下$n_{total}=m$；NA 仅用于系统/提交缺失并需单独报告）
若 $n_{total}=0$：$c_k=NA,$ $\mathrm{conf}_k=NA,$ {$\mathrm{disagree}}_k=NA$
否则：
$c_k=1$ 若$\left(n_{selected}$/$n_{total}\right)\geq0.5$ 否则 0
\[\mathrm{conf}_k=\frac{\max{\left(n_{selected},n_{total}-n_{selected}\right)}}{n_{total}}\]
\[{\mathrm{disagree}}_k=1-\mathrm{conf}_k\]
任务级聚合：
共识标签集合$C_t^{tag}=\{k$:$c_k=1\}$
整体分歧度 {$\mathrm{disagree}}_t=1-\min\{\mathrm{conf}_k$: $k\in$ $K_t\}$，取最不确定 tag 的不确定性作为保守代理
其中 $K_t=\{k$:$\mathrm{conf}_k\neq$ $NA,\$ $n_{total,k}\ge$ $n_{min}\}$，默认 $n_{min}=2$。
若 $K_t$ 为空：令 {$\mathrm{disagree}}_\mathrm{t}=NA$，并单独报告缺失率。
一致性指标 Fleiss' Kappa：
对每个 tag 独立计算 κ，报告：median κ (across tags)
期望：κ > 0.5，中等一致性（moderate agreement）左右
若某 $tag\$ $\kappa<\$ 0.4：采取预注册处理流程，避免随意切换：
(1) 频率检查：若该 tag 在 Calibration 中出现频率 < 5\%，按 per-tag 共识 $c_k=1$ 计，则并入其他低频场景，仅用于审计不用于路由分层。
(2) 若频率\ $\geq5$：不直接用该 tag 作为单独场景做细粒度路由，而是与相近 tag 合并为粗粒度场景，例如 occlusion + low_texture 合并为低可见性，并在附录中公开合并映射表。
(3) 报告与使用：该 tag 仍在 $Worker\times$ Scene 矩阵与分层表中单独报告，但标注低一致性警告；路由若引用该场景的 $r_u^{\left(s\right)}$，强制使用其置信下界 LCB 而非点估计，避免噪声场景误导。

\subsubsection{可靠度加权共识：核心干预手段}
动机：结合 Kara 的共识估计与 Liu 的工人分型思想，将可靠与不可靠工人的贡献区分开，以提升标签共识的稳健性与审计可解释性。
适用前提：默认不包含恶意标注者；通过 PreScreen 尽可能滤除随意/不稳定标注者，并在 Discussion 中承认该前提的局限。若 PreScreen 结果显示存在明显恶意或无法稳定区分可靠度，则加权共识不作为主流程决策依据，仅报告未加权多数票与一致性指标用于审计。
关键约束：
权重 $w_u$ 的来源必须独立于当前任务的标签与共识计算。
优先采用预筛选测试样本与专家参考标准得到的初始可靠度 $r_u^{\left(0\right)}$，并以此映射到权重 $w_u$；其次采用校准集的 leave-one-task-out 估计。
加权共识用于决策、审计、路由的算法输出，不作为唯一主结论的评价目标；主评价保持独立的参考标准或未加权一致性指标。
加权规则（对每个 tag k 独立）：
令 $z_{t,u,k}\in\{0,1\}$表示 worker u 是否勾选 tag k；权重 $w_u\in\left[0,1\right]$
加权投票率：$p_{t,k}^{\left(w\right)}=\frac{\sum_{u=1}^{m}{w_uz_{t,u,k}}}{\sum_{u=1}^{m}w_u}$
边界条件以避免除零且可审计：若 $\sum_{u=1}^{m}w_u=0$，则该任务在加权共识上退化为未加权多数票（或记为 NA 并单独报告；两者择一需在 Stage 1 预注册锁定）。
加权共识：$c_{t,k}^{\left(w\right)}=\mathbb{I}\left[p_{t,k}^{\left(w\right)}\geq0.5\right]$
权重映射与裁剪：预注册锁定，避免事后调参
(1) 基础映射：以 PreScreen 测试样本上的初始可靠度 $r_u^{\left(0\right)}$ 为主锚点，映射到 $w_u\in\left[0,1\right]$：
\[w_u=\min{\{}1,\max{\{}0,\left(r_u^{\left(0\right)}-0.5\right)/0.5\}\}\]
其中$r_u^{\left(0\right)}<0.5$ 视为极端不可靠，设 $w_u=0$。
(2) 上界裁剪：$w_{max}$ 由 Stage 1 PreScreen 的专家参考样本按预注册规则锁定，避免固定常数带来的任意性与事后调参空间。
候选集合：$\mathcal{W}=\{0.80,0.85,0.90,0.95,1.00\}$。
方案 A（默认，优先）：在 PreScreen 专家参考样本上执行 K 折交叉验证/交叉拟合（默认 $K=5$；样本不足时 $K=2$）。每一折仅用训练折估计 $r_u^{\left(0\right)}$与 $w_u$，再在验证折上评估加权共识与专家参考的一致性增益（$\Delta$ acc 或 $\Delta\kappa$）。
选择规则：取跨折均值最优的 $w_{max}$；若多个候选落在最优的 1-SE 范围内，取其中最小的 $w_{max}$。
方案 B：若 PreScreen 专家样本量或覆盖不足以稳定执行方案 A，则一次性随机拆分 A/B：A 半仅用于估计$r_u^{\left(0\right)}$，B 半仅用于在 $\mathcal{W}$ 上选择 $w_{max}$；拆分随机种子与清单写入预注册并公开。
锁定：最终 $w_{max}$ 在进入 Calibration/Main 前冻结，后续不再调整。
(3) 敏感性分析：附录报告 $w_{max}\in\mathcal{W}$ 全网格下结果一致性,用于稳健性披露，不再用于二次选择。
报告方式：同时报告未加权与加权的差异；并在独立参考子集上比较两者与参考标准的一致性与稳定性。
\subsection{质量与可靠度评估：LOO 原则}

可靠度评估遵循两项核心原则：其一，排除自身影响,共识计算不包含被评估者，避免 $r_u$ 系统性偏高；其二，仅用手动标注集估计,避免工具将所有人拉到同一起点。
2.4.1 任务级一致性 $\mathrm{IAA}$
对任务 t，收集所有标注者的提交结果$A_{t,u}$ 需要任务标注者数量$\geq2$，定义任务内一致性为两两 IoU 的中位数：
\[\mathrm{IAA}_t=\mathrm{median}_{u\neq v}\,\mathrm{IoU}\left(A_{t,u},A_{t,v}\right)\]
采用中位数而非均值，旨在降低极端标注的影响。$\mathrm{IAA}_t$ 高意味着任务上存在稳定共识，可用于识别专家标注者或建立参考标准。
2.4.2 共识标注和去除自身后的共识 LOO Consensus
在任务 t 上，定义共识代表为：
\[C_t=\arg\max_{u}\ \mathrm{median}_{v\neq u}\,\mathrm{IoU}\left(A_{t,u},A_{t,v}\right)\]
为了避免自我参与导致共识虚高，对每个标注者采用去除自身后的共识：
\[C_t^{\left(-u\right)}=\arg\max_{w\in U_t\setminus\{u\}}\ \mathrm{median}_{\substack{v\neq w\\ v\in U_t\setminus\{u\}}}\,\mathrm{IoU}\left(A_{t,w},A_{t,v}\right)\]
并计算：$\mathrm{IoU}_{\mathrm{LOO}}\left(t,u\right)=\mathrm{IoU}\left(A_{t,u},C_t^{\left(-u\right)}\right)$
直接与其他标注者比较：
\[\mathrm{IoU}_{\mathrm{others-med}}\left(t,u\right)=\mathrm{median}_{v\neq u}\,\mathrm{IoU}\left(A_{t,u},A_{t,v}\right)\]
$k=3$ 退化披露：当任务标注者数 $k=3$ 时，LOO 后仅剩 2 人参与共识代表选择，$\mathrm{IoU}_{\mathrm{LOO}}$ 的方差偏大且共识代表退化为二选一。为此，在 $k=3$ 的任务上同时报告 $\mathrm{IoU}_{\mathrm{others-med}}$ 与两两一致性分布作为替代一致性代理；在附录中以该替代代理重复主结论的敏感性分析，预注册锁定为一次，不引入新自由度。
2.4.3 标注者全局可靠度
标注者 u 在校准集合上的任务集合为 $T_u$
\[r_u=\mathrm{median}_{t\in T_u}\,\mathrm{IoU}\left(A_{t,u},C_t^{\left(-u\right)}\right)\]
估计条件：
仅在 $Calibration_manual$ 集合估计
仅 in-scope 样本
Bootstrap 置信区间：优先采用 BCa bootstrap，重采样次数在预注册中固定为 $B=10,000$；场景分层时采用 stratified bootstrap 保证各场景覆盖。
估计精度停机准则（预注册锁定）：除最小任务数阈值外，进一步以“可靠度估计的置信区间精度”作为 Calibration 是否追加任务的客观标准。
定义：对每个 worker u，计算 $r_u$ 的 95\% CI 半宽
\[h_u=\left(UCB\left(r_u\right)-LCB\left(r_u\right)\right)/2 。\]
停机规则：若 $h_u\le\epsilon_r$，则该 worker 在 Calibration 阶段不再追加任务；否则仅从预注册的 $Calibration_reserve$ 任务池中随机追加小批量任务，并在每次追加后重新计算 CI。$Calibration_reserve$ 是从预先固定的 $Calibration_pool$ 中划分出的保留子集，用于避免靠扩大数据池来达标。
预算上限：为避免靠无限增加数据才能达标，预注册固定每个 worker 的最大 Calibration 任务数 $N_{cal,max}$（含 reserve）。若达到 $N_{cal,max}$ 仍不达标，则停止追加并标记为 low-precision；后续路由仅使用其 LCB 或直接降级为仅分配低风险任务，并在口径 T 中披露该比例。
最小任务数阈值：N ≥ 5
冗余下限：为降低 $k=3$ 下 LOO 的退化与高方差风险，$r_u$ 与 $r_u^{\left(s\right)}$ 的主估计仅使用每任务标注者数 $k\geq4$ 的 $Calibration_manual$ 任务；$k=3$ 任务仅用于诊断性报告或附录敏感性分析，不作为主估计依据。
独立性补充：在 $PreScreen_manual$ 的测试样本子集中构建专家参考标准，得到 $r_u^{\left(0\right)}$作为初始可靠度锚点；后续的 $r_u$ 与 $w_u$ 可视为在此基础上的校准估计，从而打破可靠度与共识的循环依赖。
\subsubsection{标准 Layout 指标与门控}
为与 HoHoNet、HorizonNet、Dula-Net 等前人研究的常用指标对齐，在满足 scope = in-scope 的样本上额外报告以下标准 layout 指标：
2D IoU 与 3D IoU: HoHoNet 风格的地面投影重合度 2D 与考虑高度体积的重合度 3D
由布局几何渲染的深度图RMSE 与 $\delta$：将 layout 按几何关系渲染为全景视角下的每像素深度图，再对两份 layout 的深度图做像素级 RMSE 与阈值准确率
门控分为两类分别展示：
可计算型门控：几何归一化失败、多边形无法构造、layout rendered RMSE 生成失败等
可比型门控：点对点 RMSE 需要稳定匹配时
同时 $n_pairs_mismatch$ 不再作为标准 layout 的门控条件，不要求预测和标注角点对数一致，以避免将人工进行显著修正或工作量较大的样本排除。
结构差异相关的信息 $\mathrm{\Delta}$ $n_{pairs}$、是否存在奇数个数的点、配对失败等用作分层或诊断的变量，用于报告和反例筛选。

\subsection{结构差异信号}

定义角点对数差：
\[\mathrm{\Delta} n_{pairs}\left(t,u\right)=n_{pairs}\left(A_{t,u}\right)-n_{pairs}\left(P_t\right)\]
这是难度与纠错的参考信号，用于分层报告和反例筛选。

结构诊断代理$g_t$：
为避免未标注前无法获知 difficulty/scene的情况，本文在任何人工标注之前，仅基于模型初始化/预测结果 $P_t$ 构造结构诊断信号 $g_t$，用于首轮风险分层与触发应激模式。$g_t$ 不刻画正确性，而刻画结构性高风险，典型触发包括：多边形无法构造或自交、拓扑闭合失败、角点重复/异常、以及由 $P_t$ 导出的可计算型门控失败（如几何归一化或渲染失败）。
$g_t$ 的实现以可复现规则为主，阈值与触发集合在进入 Calibration/Main 前预注册锁定；若 $g_t$ 不可计算（如字段缺失），则记为 NA 并在口径 T 中单独报告缺失率。
\subsection{预筛选与工人画像}

在小样本条件下，少数低质量工人可能显著影响全局结论。因此本文在可靠度估计与路由之前，先通过预筛选识别并剔除不合格或不稳定的工人。
预筛选协议：
盲测任务池 $PreScreen_manual$：从 In-scope 且代表性强的样本中抽取一小批任务，要求每个候选工人完成同一批任务。
测试样本：在 $PreScreen_manual$ 中加入10–20 个带专家参考标准的样本，用于估计$r_u^{\left(0\right)}$ 并初始化 $w_u$。
$w_{max}$ 锁定协议：在候选集合 $\mathcal{W}=\{0.80,0.85,0.90,0.95,1.00\}$ 上预注册选择；默认执行方案 A（K 折交叉验证/交叉拟合）锁定 $w_{max}$，若样本不足则降级方案 B（一次性 A/B 拆分），并公开随机种子与划分清单。
测试样本设计策略：
(1) 模型预测显著错误：包括跨门扩张、拓扑崩溃、角点重复等场景，测试工人是否盲目信任初始化。
(2) 边界证据不足但仍 In-scope 的样本：门框弱或遮挡明显，测试是否随意猜测导致结构性错误。
(3) 简单样本 trivial，测试基础能力与稳定性。
专家参考标准构建：由 2 名高级标注员独立精标；若两者 $\mathrm{IoU}\geq0.9$ 则取其一致结果作为参考；若 <0.9 则由第 3 人裁决，并将该样本标注为高歧义，用于审计工人在歧义场景下的行为而非作为硬性淘汰依据。
防泄漏原则：PreScreen 与 Calibration 只用校准阶段数据设门槛与估计可靠度，不使用 Validation 或 Test 的任何信息来设阈值或选取样本。
最小门槛：预注册阈值，以便复现
门控失败率$\le\tau_{fail}$
Scope 判定一致性 $\geq\tau_{scope}$
一致性代理 $\geq\tau_r$
最小有效任务数$\geq$ $N_{pre}$
新工人准入：新工人须先完成 $PreScreen_manual$ 并通过门槛，方可进入路由池；该流程写入预注册，以避免事后选择性纳入。
工人画像 ：
基本假设与边界：系统部署在受控组织内，主要面对噪声、平庸与不稳定标注者；不以强对抗恶意攻击为默认情景。为此本文将spammer定位为可观测行为模式，如低质量/随意/系统性偏差/盲信初始化等；若 PreScreen/Calibration 显示异常模式显著，则加权共识与画像在路由中的使用将按预注册规则降级（例如仅用于审计或仅用 LCB）。

二维画像坐标：纵轴为可靠度，横轴为异常度（spammer score）。两轴均仅使用 PreScreen/Calibration 数据估计，遵循 leave-one-task-out 与时间窗协议，避免循环依赖。
(1) 可靠度轴（Reliability）：采用全局可靠度的置信下界 $LCB\left(r_u\right)$ 作为保守排序依据；$r_u$的估计、CI 与停机准则见 §2.4.3。
(2) 异常度轴（Spammer score）：定义 worker 级异常度 $S_u\in\left[0,1\right]$，由五个可审计分量组成，并在进入 Main 前预注册锁定。
$S_u=\mathrm{clip}\left(\sum_{j=1}^{5}{\alpha_jp_j\left(u\right)},\$ 0,\ 1$\right)$，默认α1=⋯=α5=0.2。
五个分量均以该 worker 在 PreScreen/Calibration 的任务集合 $T_u$ 为统计口径：
$p_{fail}\left(u\right)$：可计算型门控失败率，即其提交导致“几何归一化失败/多边形不可构造/渲染失败”等可计算型门控失败的比例。
$p_{NA}\left(u\right)$：字段缺失率，即 scope/difficulty/$model_issue$ 等关键字段出现 NA（系统缺失/漏传）的比例。
$p_{idle}\left(u\right)$：低努力代理比例。定义每任务 $active_time(t,u)$；若 $active_time(t,u)<\tau_{idle}$ 计为 1，否则 0；$\tau_{idle}$ 在 PreScreen 阶段按预注册规则锁定，例如取通过门槛工人 $active_time$ 的 10\% 分位。
$p_{trust}\left(u\right)$：盲信初始化代理比例。若 $\mathrm{IoU}_{\mathrm{edit}}\left(t,u\right)>0.95$ 且 $\mathrm{IoU}_{\mathrm{LOO}}\left(t,u\right)<0.7$，则记为高置信不改但不一致。$p_{trust}(u)$ 为该事件在 $T_u$ 中发生的比例。
$p_{var}\left(u\right)$：不稳定性代理。令 $s\left(t,u\right)=\mathrm{IoU}_{\mathrm{others-med}}\left(t,u\right)$；定义 $v_u=\mathrm{IQR}_{t\in$ $T_u}\,s\left(t,u\right)$，并令 $p_{var}\left(u\right)=\min\left\{1,\frac{v_u}{\tau_{var}}\right\}$，其中 $\tau_{var}$ 在 PreScreen 阶段按预注册规则锁定。
上述分量均为过程可观测指标，不依赖主实验的后验性能，因此可用于审计与路由而不构成信息泄露。

worker 类型（6 类，预注册规则）：
将 worker 映射为离散类型以便报告与路由约束。阈值 $\tau_{r,high},\tau_{r,low},\tau_{S,high},\tau_{S,low}$ 在 PreScreen 阶段预注册锁定。
Reliable worker：$LCB(r_u)\ge\tau_{r,high}$ 且 $S_u\le\tau_{S,low}$。
Normal worker：$LCB(r_u)$ 介于 $\tau_{r,low}$ 与 $\tau_{r,high}$，且 $S_u$ 不高。
Sloppy worker：$LCB(r_u)$ 偏低或中等，且 $S_u$ 偏高（常由 $p_{fail}$/$p_{idle}$ 拉高）。
Random spammer：$LCB(r_u)\le\tau_{r,low}$ 且 $S_u\ge\tau_{S,high}$，并伴随高波动（$p_{var}$ 高）。
Uniform spammer：$S_u\ge\tau_{S,high}$ 且表现出单一模式的系统性异常（例如长期高 $p_{trust}$ 或长期高 $p_{NA}$/$p_{fail}$）。
Partial spammer：全局 $LCB(r_u)$ 不低，但在“高风险任务桶”中显著退化。
高风险任务桶 b 仅使用标注前可得代理定义：$b\in\{g_t\$ $\mathrm{triggered}\}\times\{d_t\$ $\mathrm{high$/$low}\}$（$d_t\$ $\mathrm{high}$ 按校准集分位预注册锁定）。
对每个 b 估计 $r_u^{(b)}$ 及其 LCB，并定义 $\Delta_u=\max_b$ $r_u^{(b)}-\min_b$ $r_u^{(b)}$。若 $\Delta_u\ge\tau_{gap}$ 或在高风险桶 $LCB(r_u^{(b)})$ 明显偏低，则标记为 partial spammer。$\tau_{gap}$ 在 PreScreen 阶段预注册锁定。

与路由/共识的接口（预注册锁定）：
高风险任务（$d_t$ 高或 $g_t$ 触发）：候选池限制为 Reliable/Normal；partial spammer 在其失效的高风险桶内禁用。
低风险任务：允许 Normal/Sloppy 参与，但对高 $S_u$ 的 worker 提高抽检冗余或仅分配低风险。
Random/Uniform spammer：默认不进入主路由池；若保留则仅用于审计对照，并在共识与加权中按预注册规则置 $w_u=0$。
过程性证据：在 §4.4 固定抽样的“同任务不同画像”可视化中，对比不同类型 worker 的标注差异，验证分型与路由收益的归因链。
\subsection{场景特异可靠度（Scene-Specific Reliability）}

工人能力并非全局常数，不同场景类型下同一工人的可靠度可能存在显著差异。因此本文引入场景特异可靠度 $r_u^{\left(s\right)}$，以支持更精细的路由策略。
场景定义：以 difficulty/$model_issue$ 的共识标签为场景代理，优先采用粗粒度分层（如 occlusion、low_texture、topology_failure、overextend_adjacent），避免场景组合爆炸。保留 $S_{max}$个高频粗场景（$S_{max}$在预注册中固定，比如 $S_{max}=6$），其余低频场景并入 other；场景选择依据仅使用 Calibration 阶段的频率与一致性数据，不使用 Main 阶段结果。
估计公式：
\[r_u^{\left(s\right)}=\mathrm{median}_{t\in T_s}\,\mathrm{IoU}\left(A_{t,u},C_t^{\left(-u\right)}\right)\]
其中 $T_s$ 为场景 s 的任务集合。
估计条件：
每个场景至少 5 个任务；Bootstrap 95\% CI。
只在 $Calibration_manual$ 估计，保持与全局$r_u$ 一致。
鲁棒性保护：若某场景样本数 < 5 或 CI 过宽，则该场景下的路由退化为使用全局 $r_u$ 或更保守的 LCB，避免不稳定的 $r_u^{\left(s\right)}$ 误导策略。
主分析不引入收缩平滑超参数（如 $N_0$），不拟合 $r_u\left(d_t\right)$ 连续函数，也不使用基于 tag 相关性的层级回退。
这些扩展仅放在 讨论板块和未来工作：现有样本规模与周期不足以稳定估计，且会增加调参空间。

\subsection{任务路由策略（Scene-Aware Routing）}

本文不将自演化收敛作为主贡献，而是将路由视为固定预算下的质量控制策略：在固定预算内，将更高风险的任务优先交由更稳定的工人，并用序贯追加与停止准则把冗余花在刀刃上。核心要求是可审计、可离线模拟，且在 non-IID 与 OOD 条件下保持稳健。

为打破冷启动困境，本文将路由信号区分为“标注前可得”与“标注后可得”两类，并将二者的使用时点写入预注册：
(1) 标注前可得的风险代理（cold-start proxies）：
分布风险 $d_t$：由冻结特征 $e_t$ 与校准集分布的距离或密度估计得到，用作触发器与分层变量。
结构诊断 $g_t$：由模型初始化/预测结果 $P_t$ 的结构合法性与失败模式构造，用于识别结构性高风险。
(2) 标注后可得的人的诊断信号（post-submission diagnostics）：
$C_t^{tag}$：difficulty 与 $model_issue$ 的共识标签集合；以及基于 per-tag $conf_k$ 的任务级分歧度 $disagree_t$。
这些信号仅在任务已获得至少 2 份有效元标签后启用，用于序贯追加与解释性分层，不作为首轮派单的前置条件。
(3) 系统失败信号：门控失败类型、字段缺失率（NA 缺失单独审计）。

预注册主实现（$d_t$）：
embedding 来源：固定采用 HoHoNet 共享 pre-head 1D latent，经宽度池化与 L2 归一化后得到图像级特征；若因环境原因不可用，仅在附录给出替代 embedding 并报告其与主实现 $d_t$ 的一致性。
$d_t$ 定义：基于校准集的 kNN 距离（L2 范数）
\[d_t = \frac{1}{K}\sum_{i\in \mathrm{NN}_{K}(t)}\left\|e_t-e_i\right\|_{2}\]
kNN 超参锁定：主实现固定 $K=10$；附录做敏感性分析 $K\in\{5,10,20\}$。
敏感性分析：Mahalanobis 距离或密度估计等作为替代实现对比，放在附录或补充材料中。
严谨定位：$d_t$ 不直接判对错，不作为硬门控过滤样本；它用于审计偏移是否存在、触发应激模式、以及在 3.0 节中做分层对比。

策略设计遵循以下原则：
预筛选约束：仅在通过 PreScreen 的工人池中路由；未通过者只分配低风险任务或剔除。
保守选择：用 $r_u^{\left(s\right)}$ 的置信下界进行排序，避免小样本场景下点估计误导。
反 domination：引入负载惩罚或连续派单上限，避免少数工人被过度使用。
应激模式（stress mode）：当 $d_t$ 高、$g_t$ 触发或系统失败信号异常时进入应激模式，收紧候选工人至高置信下界子集，并提高冗余上限或触发专家复核队列。
序贯冗余与停止准则（预注册锁定）：
首轮派单：默认先派 1 份标注；若 $d_t$ 高或 $g_t$ 触发，则首轮直接提升为至少 2 份标注（确保可计算 $disagree_t$），并可同步收紧候选工人池。
追加触发：当任务累计获得$\geq2$ 份有效元标签后，计算 $disagree_t$；若 $disagree_t$ 高于阈值、出现门控失败/字段缺失（NA）异常、或有效人数不足，则每次追加 1 份标注。
停止条件：当 $disagree_t$ 低于阈值（或等价地最不确定 tag 的 $conf_k$ 达标）且 NA 缺失不触发审计异常时停止追加。
硬上限与降级：预注册固定最大冗余 $k_{max}$。若达到 $k_{max}$ 仍未达标，则不再追加，进入“专家复核/单独审计”队列，并在口径 T 中披露该类任务比例；主分析不将其静默过滤。
本文对比以下三种策略：
策略 1: Random baseline，不使用画像、场景、OOD。
策略 2: Global routing，仅用全局 $r_u$。
策略 3: PreScreen + Stratified + OOD-aware routing，路线 A 主体。

离线模拟评估协议：
数据：$Calibration_manual$/semi 每任务 $k\geq3$ 标注。
方法：模拟预筛选→序贯分配→停止，并取真实提交作为输出，不做事后挑选。
防后见之明协议：预注册锁定
Leave-one-task-out：对每个任务 t，估计 $w_u$ 与 $r_u^{\left(s\right)}$ 时不使用 t 的任何信息。
时间顺序：在数据采集阶段记录每次提交的时间戳；离线模拟时按时间构造历史窗口，仅用 t 之前的历史任务估计 $w_u$ 与 $r_u^{\left(s\right)}$ 再对 t 进行路由。
时间戳缺失处理：若时间戳字段缺失，记录为 NA 并在口径 T 中单独报告；该任务不进入时间窗口敏感性分析，但仍可进入 LOO 协议下的主分析。
冷启动：当某 worker 的历史有效任务数 <$N_{min}=5$ 时，其 $r_u$ 与 $r_u^{\left(s\right)}$ 不启用，路由退化为在通过 PreScreen 的池内随机或仅按 $r_u^{\left(0\right)}$ 粗排；当历史 $\geq$ $N_{min}$ 后再启用基于 $r_u$ 与 $r_u^{\left(s\right)}$ 的分层路由。$N_{min}$ 存在离散跳变效应，即阈值附近的标注者路由策略会发生突变；为证明结论不依赖该特定阈值，预注册固定比较 $N_{min}\in\{3,5,7\}$ 三档的敏感性分析。
指标：一致性代理、门控失败率、$active_time$ 若有、以及预算效率$\bar{k_{used}}$。
\subsection{OOS（Out-of-Scope）统计与混合判定}

对于标记为 OOS 的样本，主指标分析中予以剔除，但单独报告其比例、子类分布与判定一致性。
混合判定：
当存在 u 选 In-scope 且 v 选 OOS 时：
标记为 $scope_conflict$
人工复核协议：由两名独立高级标注员盲法复核；若分歧则由第三人裁决，并报告需要第三方裁决的比例。
报告混合判定比例

\subsection{反例筛选规则}

本文定义三类反例，各类均有明确的可复现条件：
Type 1: 低一致性
条件：$\mathrm{IAA}_t<0.7$ 且 $n_{annotations}\geq3$
可能原因：任务歧义或标注错误
Type 2: 异常编辑模式
2a 改动极大但 $\mathrm{IAA}$ 未提升：$\mathrm{IoU}_{\mathrm{edit}}<0.5$ 且 $\mathrm{IAA}_t<\mathrm{median}(\mathrm{IAA})$
2b 改动极小但与他人严重不一致：$\mathrm{IoU}_{\mathrm{edit}}>0.95$ 且 $\mathrm{IoU}_{\mathrm{LOO}}<0.7$
Type 3: 门控失败
标准 layout 指标不可计算
但图像质量正常 非证据不足
人工复核流程分为四个步骤：(1) 自动筛选符合规则的候选；(2) 专家逐一复核；(3) 确认反例类型；(4) 导出案例库并附带可视化。

\section{实验设置}

实验旨在同时验证三个研究问题：半自动初始化的效率收益（RQ1）、对质量与可靠性分布的影响（RQ2）、以及场景感知路由的收益（RQ3）。以下依次介绍数据组织、评估协议与统计检验方案。
\subsection{IID 与 non-IID 应激测试}

为检验审计与路由机制在分布偏移下的价值，本文在同一套标注预算与评估协议下同时覆盖两类部署情景。
IID baseline：训练、校准、验证来自近似同分布的数据抽样，采用常规随机抽样与分层控制。
non-IID stress test：通过数据切分显式制造难度与失效模式的分布偏移，用于检验审计与路由机制在偏移下是否更有价值。
构造方式：
为避免用标注后才产生的元标签反向决定数据切分，本文以标注前可得的代理信号构造应激情景：在验证集中提高高风险代理任务的占比，其中高风险由 $I_t^{\mathrm{OOD}}=1$ 与 $g_t$ 的触发共同定义，且 $I_t^{\mathrm{OOD}}=\mathbf{1}[d_t>\tau_d]$；阈值仅在校准阶段预注册锁定。
同时，为保证 non-IID 的语义可解释性，在标注后再用 difficulty/$model_issue$ 的共识标签作审计性验证：定义困难高风险集合 $S_{hard}$（如 occlusion、low_texture、topology_failure、overextend_adjacent 等 tag），并报告其在 IID 与 non-IID 切分中的占比差异；该验证用于解释性披露，而非用于切分本身。
$Validation_OOD$ 定义：可复现、避免硬门控，在验证集中取满足 $I_t^{\mathrm{OOD}}=1$ 的子集，主实现锁定为校准集 leave-one-out 分数的 $q=90\%$ 分位阈值；在 IID 近似条件下其预期高风险比例约为 10\%，但在真实 validation 或 non-IID 条件下其实际占比允许偏离 10\%。
IID：保证训练与验证中由 $d_t$/$g_t$ 定义的高风险代理占比接近，或在预先声明的误差范围内；并报告事后审计得到的 $S_{hard}$ 占比作为一致性检查。
non-IID：在验证集中人为提高高风险代理任务（高 $d_t$ 或 $g_t$ 触发）的比例，使其在该验证集内占比显著偏高；训练集与校准集则保持较低的高风险代理占比。随后用标注后的 $S_{hard}$ 占比与失败模式分布作审计性验证，说明该切分确实带来更高的困难/失效暴露。
核心比较：
OOD-aware 的审计任务：报告 $d_t$ 的分布，并给出 $d_t$ 与门控失败、反例率、$\mathrm{IAA}_t$ 的相关性，相关而非因果，用于证明分布偏移的存在，以及偏移条件下系统更容易失效。
防 selection bias：无论 $d_t$ 高低，样本仍进入口径 T/I/M 的透明披露；OOD 仅用于分层报告与策略触发，不用于静默过滤。
3.1 条件分组与数据组织（$\mathbit{N}_{\mathbit{pool}}=\mathbf{458}$）

记号：$N_{img}$ 为唯一图像数，$N_{task}$ 为任务实例数，k 为每图标注份数，$\left|U\right|$ 为 PreScreen 候选工人数。
主方案分组：
Stage 0 Pilot：$N_{img}=15$，$k=2$；仅用于流程验证，不进入主结论。
Stage 1 PreScreen：$N_{img}=50$（含 $Trap=15$）；用于 $r_u^{\left(0\right)}$ 锚点、工人筛选与门槛锁定。
Stage 2 Calibration：$N_{img}=30$，manual/semi 同图各 $k=4$；用于 $r_u$ 与 $r_u^{\left(s\right)}$ 估计及 RQ2 对照（避免 $k=3$ 下 LOO 退化导致的高方差估计）。
Stage 3 Main：
RQ1 Test：$Manual_Test$ 与 $SemiAuto_Test$ 同图 $N_{img}=100$，$k=1$。
RQ3 Validation：$Validation_{semi}$ ($N_{img}=60,$ $k=3)$；$Gold_{manual}$ ($N_{img}=20,$ $k=3)$ 用于稳健性对照。
预注册约束：
Trap 与 Main 物理隔离；Calibration/Validation/Test 图像不重叠；Pilot/PreScreen 不进入 Main 对比。
Hard 子集 H 从 Validation 构造，要求$\left|H\right|\geq30$；入选依据为高 $d_t$ 或高 $g_t$，并在标注后验证 $S_{hard}$ 占比$\geq70$。
OOD 仅用于分层与触发，不用于静默过滤；全部样本进入 T/I/M 披露。
规模说明：核心主分析唯一图像约 190（Calibration 30 + Validation 60 + Test 100），其余用于前置锚点、流程验证与稳健性。
数据与评测口径锁定：
test 与 $test_no_occ$ 的 img 集合一致，差异仅在 $label_cor$ 修正。
标准 benchmark 主报告使用 test/$label_cor$；$test_no_occ$/$label_cor$ 作为附录敏感性分析。
3.2 参与者与 worker mix 控制

培训内容包括 Label Studio 基础操作、仅标注相机所在房间的规则、以及 OOS 分流规则。
分配与记录：统一培训后不再引入新手与熟手分层。为控制 worker mix 与能力差异混杂，采用 worker 级别的平衡分配：每位通过 PreScreen 的 worker 在 Manual 与 Semi 两种条件下分配到近似相同数量的任务，并随机化任务顺序。同时记录每个 worker 的任务数、完成时间与返工率，并报告每批 worker 组成以供 worker mix 控制。
新工人准入：先完成 $PreScreen_manual$ 并通过门槛，方可进入 Calibration/Validation 的路由池；该流程写入预注册，以避免事后选择性纳入。
3.3 预注册协议（Primary Endpoints）

3.3.1 RQ1 主终点 效率
指标：active time 中位数差异 Manual vs Semi
统计：Mann-Whitney U test 非参
效应量：Hodges–Lehmann 中位数差
CI：95\% CI，优先 BCa bootstrap，重采样次数预注册固定
worker mix 控制：通过 3.2 的 worker 级别平衡分配，减少条件间能力差异带来的混杂；结果中报告两条件下 worker 构成与任务覆盖的平衡性。
风险控制：回归分析 time ~ $\mathrm{IoU}_{\mathrm{edit}}$，检验时间缩短并非源于随意修改
交叉验证：一致性验证,开展自定义计时与原生计时的一致性校验，计算自定义统计的 $active_time$ 与 Label Studio 原生 $lead_time$ 的 Spearman 相关系数，同时分析二者比值的中位数及四分位距（IQR），以此验证自定义活跃标注时间的采集逻辑无系统性偏离。本校验实施成本低，且可有效排除计时采集环节的系统性实现误差。

3.3.2 RQ2 主终点 可靠性
指标：$\mathrm{IAA}_t$ 分布差异 $Calibration_manual$ vs $Calibration_semi$
统计：Kolmogorov–Smirnov test
效应量：$\Delta\,\mathrm{median}(\mathrm{IAA})$
反例数量：三类反例比例差异 Chi-square test
RQ2b：加权共识增益，作为干预而非评价目标
评估集合：PreScreen 的测试样本（含专家参考标准）及其扩展复核子集。$w_{max}$的锁定与性能评估采用预注册的方案 A（样本不足时降级方案 B）；在每个验证折（或 B 半）上，ru0与 $w_u$ 的估计均不使用该验证部分标签，保证评估独立性。
主终点：加权与未加权共识对专家参考标准的一致性提升，如 per-tag accuracy 或 Cohen's $\kappa$ 的提升。
统计：按任务重采样的 BCa bootstrap 得到 $\Delta\kappa$ 或$\Delta$ acc 的 95\% CI；若 CI 不跨 0 且效应量 $\geq$ $MDE_\kappa$，默认 0.05，则判为工程上有意义的提升。
反例检出：对 2.10 的反例规则产生的候选集，比较加权与未加权的检出率差异，使用 McNemar's test 或 bootstrap CI。
3.3.3 RQ3 主终点 路由
指标 1 主：一致性代理提升 Stratified vs Random: $\mathrm{IoU}\left(A_{t,u},$ $C_t^{(-\hat{u})}\right)$
统计 1：Paired t-test 每任务配对；仅在 $k\geq3$ 任务上计算
效应量：Cohen's d 标准化的均值差效应量
指标 1b：预算效率，序贯停止，平均每任务使用标注份数 $\bar{k_{used}}$ 或在固定预算下可覆盖的任务数
指标 2：门控失败率降低
统计 2：McNemar's test 麦克尼马尔检验 配对二元
指标 3：Active time 变化，若有时间数据
应激测试报告：non-IID
在3.0 定义的 $Validation_OOD$ 子集上重复上述对比，以失败率与反例拦截率作为更直接的系统收益证据。
3.4 多重检验处理

策略：预先限制检验集合，避免 FDR 复杂度
RQ1：1 个主检验 Manual vs Semi = 1 个
RQ2：1 个主检验 IAA 分布 + 3 个反例类型检验 = 4 个
RQ3：3 个主检验 一致性代理 / Failure/Time = 3 个
总计：8 个计划检验
校正：Bonferroni α = 0.05/8 = 0.00625 保守
\subsection{效应量门槛与 MDE}

最小可解释差异 MDE，示例阈值，以试运行与先验为准：
Time: $\geq8$ 相对改进
一致性代理: $\geq0.05$
Failure rate: $\geq5\\$% 绝对下降
报告原则：只有同时满足 p<0.05 且效应量超过 MDE 时，才声称显著改进。
MDE 的确定采用分级锁定流程：
Stage 0 Pilot：预研，$N\approx20-30$，估计工具可用性与噪声底线，并给出候选阈值 $\tau_{0}$，如效率相对提升、失败率绝对下降等。
Stage 1 PreScreen：测试样本校准，独立锚点，在测试样本上估计工人表现方差与可达上限，并将候选$\tau_0$ 调整并锁定为最终阈值 $\tau_{final}$，即 MDE。同时在此阶段锁定 $w_u$的映射规则，并按 §2.3.2 的预注册流程锁定裁剪阈值 $w_{max}$，避免事后调参。
Stage 2 Main Experiment：正式实验，严格使用 $\tau_{final}$ 作为工程上值得的判据，且 PreScreen 与测试样本不得进入 RQ3 的路由效果评估集合，防 data leakage。
实施说明：三段式流程允许软件未完全自动化；可通过人工分发不同任务包完成阶段切换，但必须保证不手动修改标注内容，且阶段边界与样本隔离可审计。

\section{报告与可审计输出}

本文的报告体系以可审计性为核心，通过三级口径、分层矩阵与反例库构成完整的证据链。
\subsection{三级口径 T/I/M}

为避免因指标不可计算而静默丢弃样本，本文采用三级口径透明披露所有数据：
口径 T（Total）——透明披露：包含所有收集到的标注，不因指标无法计算而静默丢弃。报告 OOS 比例、门控失败原因分布、缺失率。
口径 I（In-scope）：仅包含 $scope=in-scope$ 且 scope 非缺失的样本，用于主指标分析（效率、IAA、$r_{u}$）。OOS 样本单独报告比例与子类。
口径 M（Model-usable）：在口径 I 基础上排除可计算型门控失败，用于标准 layout 指标（2D/3D IoU）。可比型门控失败仅影响对应指标。
$Active_time$ 数据质量审计：在口径 T 中以审计表形式报告 $active_time$ 日志的异常记录比例，包括 $task_id$ = unknown、$annotator_id$ = unknown、$project_id$ = unknown、$script_version$ 缺失、以及多 session (task, annotator) 对的占比，使数据质量可追溯而非仅口头声称。
\subsection{分层报告表}

本文通过四张分层表多角度展示结果：
表 A：按 Δ$n_{pairs}$ 分层——覆盖率、门控失败原因、改动幅度统计。
表 B：按 difficulty/$model_issue$ 分层——承认标签的主观性（每人主观感受不同），并报告 Fleiss' Kappa 验证一致性。
表 B2：混淆/冲突矩阵（低成本审计项）
Scope 冲突矩阵：In-scope vs OOS 及 OOS 子类冲突率
Tag 冲突统计：高频 difficulty/$model_issue$ 的共现与冲突比例
用途：定位高不确定性来源，支撑专家复核优先级；不作为硬门控
表 C：Worker $\times$ Scene 可靠度矩阵
行：$worker_id$
列：场景 occlusion, $low_texture,$ ...
单元格：$r_{u}^{(s)}$ $\pm$ 95\% CI
标注：N < 5 的格子标记为 --
类型标注：在每个 worker 行给出画像类型（Reliable/Normal/Sloppy/Partial Spammer/Random Spammer/Uniform Spammer）及两轴坐标（$LCB(r_u),$ $S_u$）；并在矩阵中标注其在不同场景/风险桶下的局部失效。
局部失效操作定义：预注册，满足以下任一条件即标记为 $failure_{u,s}=1$：
(1) 绝对阈值：$r_{u}^{(s)}$ < 0.70；
(2) 相对下降：$r_{u}^{(s)}$ < $r_{u}$ - 0.10；
(3) 保守显著性：$r_{u}^{(s)}$ 的 95\% CI 上界 < $r_{u}$ 的 95\% CI 下界。
图 D：工人画像散点图（LabelInspect 风格）
横轴：异常度 $S_u$（spammer score）
纵轴：可靠度 $LCB(r_u)$
颜色/形状：6 类 worker types
\subsection{反例案例库（Anomaly Case Library）}
输出内容：
反例清单 CSV：$task_id,$ type, $\mathrm{IAA}_{t},$ $\mathrm{IoU}_{edit},$ $consensus_tags,$ 原因
可视化：每个反例的标注对比图 3-4 个 worker 叠加
统计：三类反例的频率分布
人工复核笔记：专家确认的真正原因
作用：必须说明其因果与证据链地位
(1) 可审计筛选：依据 2.10 的可复现规则产生候选集，不依赖专家直觉选取样本。
(2) 专家复核：逐条确认是否为真正异常，并记录可解释原因与归类标签。
(3) 典型失败模式总结：形成可检索的失败模式词表，如跨门扩张、拓扑崩溃、配对失败、角点漂移等，并统计其在 IID 与 non-IID/OOD 分层下的频率变化。
(4) 反馈闭环：不等同于重训练闭环，将失败模式反馈到标注指南、培训素材、路由规则，例如在某类失败高发场景提高冗余或限制工人池。
(5) 过程性证据：在反例库中固定输出同一任务、不同画像类型工人的对比可视化，用来解释为何预筛选与分型是必要的。

\subsection{过程性证据与归因链验证}
目的：将工人画像（$LCB(r_u)$ vs $S_u$ 与 6 类分型）、Worker×Scene 矩阵、反例库与路由收益串联为可审计的证据链，而非仅报告最终 p 值。
输出：
同任务不同画像对比：选取若干典型任务，包含 OOD 与高风险，展示 reliable vs sloppy 等类型的标注差异。
案例抽样规则：预注册，避免选择性展示，固定抽取 15 个任务用于正文展示，5 个高 $\mathrm{IAA}_{t}$ 共识明确，5 个低 $\mathrm{IAA}_{t}$ 分歧严重，5 个高 $d_{t}$ OOD 风险高；每个任务展示 3–4 名来自不同画像类型的标注叠加对比。
失败模式溯源：对高频反例类型，回溯到对应场景与工人类型，解释为何会发生、如何被预筛选或路由缓解。
路由收益归因：证明 stratified routing 的收益主要来自高风险任务分配给更稳定工人与 stress mode 增冗余或复核，而不是静默过滤。

\section{标注数据重训练}
定位为探索性分析和未来工作，不作为主贡献，放到 Discussion 或 Appendix。
\subsection{目标}
如果进行重训练，主要目标包括：利用反例案例库筛选训练数据，降低结构性错误如跨门扩张和配对失败，降低门控失败率。
\subsection{评估方式}
采用固定的 Test set 进行评估以避免 leakage，对比旧模型 $M_{0}$ 和新模型 $M_{1}$，报告 $model_issue$ 比例、$\mathrm{IoU}_{edit}$ 改善幅度以及门控失败率变化。
\subsection{当前不做的原因}
尽管候选池已扩展到 458 张，但本研究的主目标仍是首先将 RQ1–RQ3 的规程、审计与路由收益构建为可复现的证据链。若引入重训练，会显著增加方案自由度（训练数据选择、训练超参数、早停与模型选择等），并带来更复杂的 leakage 风险控制与工程成本，不利于在既定时间线内完成主实验。
未来工作：在更大规模（例如 $\ge$ 800 张）且跨数据源的样本上，结合反例库与审计信号开展重训练/主动学习闭环，并保持与本文一致的预注册与隔离协议。

\section{讨论与局限性}

本章预先回应审稿人可能提出的质疑，对研究设计中的关键假设、方法边界与已知局限进行自我审查。
\subsection{Meta-Label 的主观性}

Difficulty 与 $Model_issue$ 依赖标注者主观判断，这是本研究的内在局限。
缓解措施：用 Fleiss' Kappa 验证一致性，期望 $\kappa$ > 0.5 表示中等一致性；统一培训并提供案例示范；通过 Per-tag 共识降低个体偏差。
本文的立场是：即使存在测量噪音，这些标签作为风险信号和场景代理仍具备实用价值。
\subsection{样本量限制}
当前可用候选池为 458 张（$mp3d_layout$/test）。在 10 人与紧周期约束下，主方案的核心对照集合使用约 190 张唯一图像（$Test=100$，$Calibration_pool=30$，$Validation=60$；另含 $PreScreen=50$ 与 $Pilot=15$ 用于前置锚点与流程验证）。$Calibration_pool$ 内部再划分出 reserve 子集，用于“估计精度停机准则”下对少数 CI 过宽的 worker 追加小批量任务；该追加不引入新的唯一图像池。多标注主要集中在 $Calibration_manual$/semi（$k=4$）与 $Validation_semi$（$k=3$），其中 $r_{u}$ 与少数高频粗场景 $r_{u}^{(s)}$ 的主估计仅使用 $Calibration_manual$（$k\geq4$），$Validation_semi$ 主要用于离线路由模拟与应激情景对比。
影响体现在：某些低频场景的 $r_{u}^{(s)}$ 仍可能因为样本数 < 5 或 CI 过宽而不可用；因此路由在该场景下退化为使用全局 $r_{u}$ 或其置信下界 LCB，并透明报告退化比例。
补充说明：为避免在小样本下引入额外可调超参数而造成“软门控调参”质疑，主分析采用保守的 LCB/退化策略，而不采用需额外超参数的收缩平滑或基于相关性的层级回退；相关扩展仅在未来工作讨论。
透明披露：在 Results 中报告每个场景的样本量和置信区间宽度。
\subsection{离线路由模拟的局限}
方法是在多标注数据上进行模拟分配后取该 worker 真实标注的离线路由评估。
局限是不属于真实在线路由，worker 行为可能因任务分配而改变。
辩护理由：离线模拟是路由研究的标准方法，参考 Kara et al. 和 Liu et al. 的工作；避免重新标注的巨大成本；为未来在线实验提供先验。离线模拟回答的是"在固定已采集行为数据下，策略的可解释收益下界"；路由对标注者行为的正向在线效应（如激励与适应性改变）属于未来工作范畴，不混入主结论。
\subsection{Convex Hull 的局限性}
问题在于凸包外包凹形，降低对跨门扩张的敏感度。
解决办法：明确定位 $\mathrm{IoU}_{edit}$ 为辅助指标；主结论依赖 $\mathrm{BoundaryRMSE}$ 与标准 layout 指标；在反例筛选中采用多指标组合而非仅依赖 IoU。

\subsection{基本假设与让步}
受控组织假设：系统部署在受控组织内，如企业内部标注团队，主要面对噪声、平庸、不稳定而非强对抗的恶意攻击者。
让步1：加权共识的局限，在极端 OOD 场景下可能出现多数标注者同时犯错而专家意见被淹没的情况；因此本文将加权共识定位为干预手段，并在 $d_{t}$ 高时通过 stress mode 提升冗余或触发专家复核来缓解。
让步2：循环论证风险的规避，通过 PreScreen 的测试样本与专家参考标准获得 $r_{u}^{(0)}$ 与 $w_{u}$ 的独立锚点，并采用 leave-one-task-out 的估计协议，避免将同一任务的标签既作为推断依据又作为评价结论。
让步3：本文不在主分析中估计“工人 OOD 韧性曲线”这类需要在高 $d_{t}$ 区间大量覆盖的数据驱动函数（例如 $r_{u}(d_{t})$ 的衰减斜率），仅将 $d_{t}$ 作为触发器与分层变量；更精细的韧性建模留待更大规模与更长周期的后续研究。

\section{与审稿员预期批评的预设回答}

以下针对审稿人可能提出的六个典型问题逐一给出回应。
问题1：Meta-label 会显著增加工作量吗？
回答：不会。Scope 是必要决策字段；Difficulty 与 $Model_issue$ 属于修改时的同步记录。额外时间预计 <10 秒，占平均标注时长约 8\%。
问题2：多选标签如何定义共识？
回答：采用 Per-tag 二元共识。每个 tag 独立多数票，报告 per-tag Fleiss' Kappa。NA 不计入有效人数并单独披露缺失率。
问题3：样本量有限，结论是否可靠？
回答：主终点预注册，使用 BCa bootstrap CI，报告效应量与 MDE，不以 p-value 作为唯一依据；同时披露各场景样本量与 CI 宽度。
问题4：离线路由模拟是否弱于在线实验？
回答：离线路由是当前阶段可复现且成本可控的标准做法，可用于形成在线实验先验；本文明确其边界并保留后续在线验证。
问题5：加权共识是否循环论证，OOD 下专家会否被淹没？
回答：权重锚点来自 PreScreen 测试样本与专家参考标准，并采用 leave-one-task-out，避免同任务自证；$w_{max}$ 在 PreScreen 阶段按预注册的方案 A（不足则方案 B）锁定后冻结，避免事后调参；高 $d_{t}$ 场景进入 stress mode，提高冗余或触发专家复核。
问题6：改进幅度是否过小？
回答：以预注册 MDE 判定工程意义；只有效应量跨过 MDE 且统计证据一致时，才声称改进成立。

\section{预期贡献总结}

本文的预期贡献涵盖规程设计、系统设计、方法学与实证发现四个层面。
贡献 0：规程设计。提出Pilot→PreScreen→Main 的三段式可审计规程，在进入主实验前预注册并锁定阈值、超参数与权重映射，允许流程在软件未完全自动化情况下通过任务包分发执行，但必须保留可审计的样本隔离与分发清单。
贡献 1：系统设计。提出可审计的半自动标注体系。三级报告口径 T/I/M 避免 selection bias；所有门控失败和 OOS 可追溯；Meta-label 作为决策副产品提供分布监控信号。
贡献 2：方法学。预筛选测试样本、场景特异可靠度、OOD触发路由的组合。$r_{u}^{(s)}$ 的估计协议与 LOO 原则；Per-tag 共识与分歧度计算；OOD 风险分数 $d_{t}$ 的预注册主实现与敏感性分析框架。
贡献 3：实证发现。在全景布局标注任务上验证，以最终实证结果为准。半自动初始化在相同图像对照下可降低 active time；场景与分布感知路由在相同预算下可提升一致性代理并降低失败率；反例案例库揭示三类主要失效模式，并形成过程性证据链。
对标文献：Kara et al.（JAIR 2018）聚焦主动学习与共识估计，本文在此基础上引入场景分层；Liu et al.（InfoVis 2018）聚焦 worker 分型与交互验证，本文采用离线路由模拟加以拓展。创新点在于将可审计性、场景感知与元标签共识进行系统性整合。


\end{document}
